% =====================================================================
%  CARNET D'ENTRAÎNEMENT — FICHE 6 : Les limites
%  Thème 1 — CORRIGÉ — Niveau FACILE
% =====================================================================
\documentclass[11pt,a4paper]{article}
\usepackage[utf8]{inputenc}
\usepackage[T1]{fontenc}
\usepackage{lmodern}
\usepackage[french]{babel}
\usepackage{amsmath,amssymb,mathtools}
\usepackage{icomma}
\usepackage{xcolor}
\usepackage{tikz}
\usepackage{pgfplots}
\usepackage[most]{tcolorbox}
\usepackage{enumitem}
\usepackage{array}
\usepackage{booktabs}
\usepackage{fancyhdr}
\usepackage[hidelinks]{hyperref}
\usetikzlibrary{arrows.meta,calc,intersections}
\pgfplotsset{compat=1.18}
\usepackage[a4paper,margin=2.2cm,headheight=26pt,headsep=10pt]{geometry}

\definecolor{primary}{RGB}{223,87,99}
\definecolor{primarydark}{RGB}{150,42,55}
\definecolor{excol}{RGB}{0,135,90}
\definecolor{curvecol}{RGB}{40,80,190}

\tcbset{boxbase/.style={enhanced,breakable,boxrule=0.9pt,arc=2.5pt,
   left=3mm,right=3mm,top=2.3mm,bottom=2.3mm,
   fonttitle=\bfseries,coltitle=white,toptitle=0.6mm,bottomtitle=0.6mm}}
\newtcolorbox[auto counter]{correction}[1][]{boxbase,colback=excol!5,colframe=excol,
  title={\textsc{Corrigé}~\thetcbcounter\ifx\relax#1\relax\relax\else\textmd{ --- \itshape #1}\fi}}

\pagestyle{fancy}
\fancyhf{}
\fancyhead[L]{\small\textcolor{primarydark}{\textbf{Fiche 6} — Les limites}}
\fancyhead[R]{\small\textcolor{primarydark}{\textsc{Thème 1 — Corrigé Facile}}}
\fancyfoot[C]{\small\thepage}
\renewcommand{\headrulewidth}{0.8pt}
\renewcommand{\headrule}{\hbox to\headwidth{\color{primary}\leaders\hrule height \headrulewidth\hfill}}

\newcommand{\R}{\mathbb{R}}

\begin{document}

\begin{center}
{\Large\bfseries\color{primarydark}Thème 1 — Notations, lecture \& calcul direct}\\[2pt]
{\large\color{excol}Corrigé — Niveau Facile}
\end{center}
\vspace{6pt}

\begin{correction}[calcul direct par remplacement]
On remplace directement $x$ par la cible (les fonctions sont définies au point).
\begin{enumerate}[label=\alph*),leftmargin=2em,itemsep=3pt]
  \item $\displaystyle\lim_{x\to 2}(3x^2-x+4)=3\cdot 4-2+4=12-2+4=\boxed{14}$.
  \item $\displaystyle\lim_{x\to -1}\frac{2x+5}{x^2+3}=\frac{2(-1)+5}{(-1)^2+3}=\frac{3}{4}=\boxed{\tfrac34}$
        \ (le dénominateur vaut $4\neq 0$).
  \item $\displaystyle\lim_{x\to 4}(\sqrt{x}+x)=\sqrt{4}+4=2+4=\boxed{6}$.
  \item $\displaystyle\lim_{x\to 0}(x^3+2x-7)=0+0-7=\boxed{-7}$.
\end{enumerate}
\end{correction}

\begin{correction}[lire et écrire une notation]
\begin{enumerate}[label=\alph*),leftmargin=2em,itemsep=3pt]
  \item « Lorsque $x$ devient de plus en plus grand (tend vers $+\infty$), $f(x)$ se rapproche de $3$. »
        (Géométriquement : la droite $y=3$ est asymptote horizontale.)
  \item « Lorsque $x$ se rapproche de $2$ \emph{par la gauche} (par valeurs inférieures à $2$), $f(x)$
        devient de plus en plus grand en valeur négative, c.-à-d.\ tend vers $-\infty$. »
  \item $\displaystyle\lim_{x\to 5^{+}} g(x)=0$.
\end{enumerate}
\end{correction}

\begin{correction}[limites latérales d'une fonction par morceaux]
\begin{enumerate}[label=\alph*),leftmargin=2em,itemsep=3pt]
  \item À gauche de $2$ on a $x<2$, donc $f(x)=x+1$ : $\displaystyle\lim_{x\to 2^{-}}f(x)=2+1=3$.
  \item À droite (et en $2$) on a $x\geq 2$, donc $f(x)=3x-2$ : $\displaystyle\lim_{x\to 2^{+}}f(x)=3\cdot 2-2=4$.
  \item Les deux limites latérales sont \emph{différentes} ($3\neq 4$) : \textbf{$f$ n'a pas de limite
        (globale) en $2$}. (La courbe présente un saut de hauteur $1$.)
\end{enumerate}
\end{correction}

\begin{correction}[lecture graphique]
La courbe se rapproche de la droite $y=2$ aux deux bouts, et passe par l'origine.
\begin{enumerate}[label=\alph*),leftmargin=2em,itemsep=3pt]
  \item $\displaystyle\lim_{x\to +\infty} f(x)=2$.
  \item $\displaystyle\lim_{x\to -\infty} f(x)=2$.
  \item $\displaystyle\lim_{x\to 0} f(x)=0$ \ (la courbe passe par le point $(0,0)$, sans trou ni saut).
\end{enumerate}
\end{correction}

\begin{correction}[vérifier une continuité]
$g(x)=\dfrac{x^2-1}{x+3}$.
\begin{enumerate}[label=\alph*),leftmargin=2em,itemsep=3pt]
  \item En $x_0=1$, le dénominateur vaut $1+3=4\neq 0$ : donc $1\in D_g$ et $g$ est définie en $1$.
  \item $\displaystyle\lim_{x\to 1} g(x)=\frac{1^2-1}{1+3}=\frac{0}{4}=0$, \quad et \quad
        $g(1)=\dfrac{1-1}{1+3}=0$.
  \item La limite ($0$) est égale à $g(1)$ ($0$) : \textbf{$g$ est continue en $1$}. (C'est le cas général
        d'une fonction rationnelle en un point où le dénominateur ne s'annule pas.)
\end{enumerate}
\end{correction}

\end{document}
